% V2 derivation. This is a leading-order conditional criterion, not a theorem
% for the full RF equations or a stability result.
\section{Leading radius-eliminated criterion}
Let $\tau=\Omega_1t$, $r=\rho/R$, $\alpha_i=(q_i/m_i)/(q_1/m_1)$,
and normalize the leading $k=+2,-4$ MFS coefficients to one. Write
\begin{equation}
K_{1i}=1+\alpha_i\beta,\quad K_{2i}=\nu^2+\alpha_i\beta,\quad
D_{1i}=K_{1i}^2+\lambda_i^2,\quad
D_{2i}=K_{2i}^2+(\lambda_i\nu)^2.
\end{equation}
At zero slow angular speed, with Coulomb, radial drag, centrifugal force,
higher harmonics, and finite-size effects omitted, the mean radial balance is
\begin{equation}
0=\alpha_i\beta r_i-
\frac{4\alpha_i^2A^2K_{1i}}{D_{1i}}r_i-
\frac{48\alpha_i^2B^2K_{2i}}{D_{2i}}r_i^5.
\end{equation}
For general leading coefficients $c_2,c_{-4}$, the factors $A^2$ and $B^2$
are replaced by $A^2|c_2|^2$ and $B^2|c_{-4}|^2$. Define
\begin{equation}
U_i=\alpha_i\beta-\frac{4\alpha_i^2A^2|c_2|^2K_{1i}}{D_{1i}},\qquad
V_i=\frac{48\alpha_i^2B^2|c_{-4}|^2K_{2i}}{D_{2i}}.
\end{equation}
A positive algebraic radius exists iff $U_i>0$ and $V_i>0$; it is
\begin{equation}
r_i^4=\frac{U_i}{V_i}=
\frac{\left(\alpha_i\beta D_{1i}-4\alpha_i^2A^2|c_2|^2K_{1i}\right)D_{2i}}
{48\alpha_i^2B^2|c_{-4}|^2K_{2i}D_{1i}}.
\end{equation}
The physical aperture additionally requires $r_i+a_i/R<r_{\rm ap}/R$.
Small-micromotion validity is a separate condition, e.g. a posteriori
$\max|\boldsymbol\xi_{\rm micro}|/r_i\ll1$ and no unresolved near-resonant
combination modes. A positive root alone is not a stability proof.

The linearized scalar radial force at fixed $\omega=0$ has slope
$F_r'(r_i^*)=-4U_i<0$, a restoring slope in that restricted radial section.
The full slow dynamics couples radius, radial velocity, angular momentum, and
the Fourier response; stability must use its coupled Jacobian. In particular,
eliminating $r_i$ from a stationary algebraic balance does not eliminate a
radial dynamical degree of freedom.

For the leading traveling modes, the slow-orbit RF torque at zero mean orbit
speed is proportional to
\begin{equation}
\mathcal T_i(0)=\frac{4\alpha_i^2A^2|c_2|^2\lambda_i}{D_{1i}}r_i^2-
\frac{48\alpha_i^2B^2|c_{-4}|^2\lambda_i\nu}{D_{2i}}r_i^6.
\end{equation}
Its zero is
\begin{equation}
r_{\rm cr,i}^4=\frac{A^2|c_2|^2D_{2i}}
{12B^2|c_{-4}|^2\nu D_{1i}}.
\end{equation}
Thus $S_i=(r_i/r_{\rm cr,i})^4$ obeys the exactly simplified identity
within this leading finite-$\lambda$, finite-DC radial/torque model:
\begin{equation}
S_i=\frac{\nu}{K_{2i}}
\left[\frac{\beta D_{1i}}{4\alpha_i A^2|c_2|^2}-K_{1i}\right].
\end{equation}
The $B$ and $c_{-4}$ factors cancel because the radial $k=-4$ force and its
zero-speed slow torque both scale as $B^2|c_{-4}|^2r^6$. This algebraic
cancellation does not survive as a statement that $B$ is physically
irrelevant: $B$ sets $r_i$, aperture margin, and the validity of the orbit
ansatz. It is also singular at $B=0$, where the confining $r^5$ radial term
vanishes and the finite-radius root may disappear.

If $c_2=c_{-4}=1$, $\lambda_i\to0$, and $|\alpha_i\beta|\ll1$ so that
$K_{1i}\simeq1$, $K_{2i}\simeq\nu^2$, $D_{1i}\simeq1$,
$D_{2i}\simeq\nu^4$, then
\begin{equation}
S_i\simeq\frac{\Pi/\alpha_i-1}{\nu},\qquad
\Pi=\frac{\beta}{4A^2}.
\end{equation}
For $\alpha_1=1$, $0<\alpha_2<1$, and $\nu>0$, the strict M0 sign
conditions $0<S_1<1$ and $S_2>1$ reduce to
\begin{equation}
\max\{1,\alpha_2(1+\nu)\}<\Pi<1+\nu.
\end{equation}
This predicts the signs of the zero-speed slow torque. Transfer to stationary
orbit signs $\omega_1>0,\omega_2<0$ additionally requires that the total
slow torque have a relevant root and positive effective angular damping
$D_i^{\rm eff}=-\partial_\omega\mathcal T_i|_{\rm root}>0$, with no other
nearby root or non-monotone branch. It is not an unconditional sign theorem
for finite-speed M3 or the full RF system.
